\scanpage{44}

\diagram{\begin{tikzcd}[ampersand replacement=\&]
\& X\arrow[d]\\
S\& T\arrow[l,"f"']
\end{tikzcd}}

$f_*(X/T)$ \qquad $T$ fini sur $S$

\diagram{\[
\begin{array}{llll}
\text{si}&X/T\ \text{est affine},&f_*(X/T)\ \text{aussi}.&\\
\text{si}&\text{---}\left\{\begin{array}{l}\text{projectif}\\\text{propre}\end{array}\right.\ast&
\left.\begin{array}{l}\text{---}\\[-.4ex]\phantom{\text{---}}\end{array}\right\}&\text{\uline{pourvu que} } f\ \text{\uline{séparable}}\ (*)\\
\text{si}&\text{quasi-projectif},&\text{---}&
\end{array}
\]}

% slanted
\margin{donner un énoncé ``général'' dans le cas où $f$ est fini \uline{étale}}

\unsure{Le calcul|Le foncteur|Notons} $f_*$ est exact à gauche, \deleted{mais pas en tout exact} \inserted{en particulier}

\[
f_*(T/T)=S,\quad f_*(X\times_T Y/T)=f_*(X/T)\times_S f_*(Y/T),\quad f_*(X\times_Z Y/T)=f_*(X/T)\times_{f_*(Z/T)}f_*(Y/T)
\]

Le foncteur $f_*$ transforme imm. ouvertes en imm. \deleted{\unsure{fermées}} \inserted{ouvertes} et \unsure{sous} une légère réserve, [\deleted{\unsure{pour que|il ne suffit}} $f_*(X/T)=\bigcup_i f_*(U_i/T)$, où les $U_i$ \inserted{affines} recouvrent $X$] imm. fermées en imm. fermées. \deleted{\unsure{(} \unsure{d'ailleurs, à réserve,|et, à réserve,}}

\margin{$*$ Se voit par descente : Si $Y/S$ est tel que $f^*(Y)/S'$ soit propre resp. projectif, alors $Y/S$ l'est [Dans le deuxième cas, on suppose $T/S$ fini et plat sur $S$ \uline{loc. noeth} \deleted{\unsure{(}}], (i) \deleted{\unsure{fau|faut}} devrait dans le cas actuel pouvoir l'exorciser, en se ramenant au projectif type, \deleted{\illegible{2}} travaillant sur $\underline{\mathbb Z}$ en définitive ; il faut prouver ceci : Si un faisceau loc. libre inversible $\underline L$ sur $X$ est $T$-ample, alors $f_*(\underline L)$ (qui est un faisceau loc. libre de rang $d$ sur $f_*(X/T)$ est tel que $\det f_*(\underline L)$ est $S$-ample \ldots}

\margin{Dans le premier cas, on utilise le lemme de Chow, en notant que si $X'\longrightarrow X$ est surjectif, alors $f_*(X')\to f_*(X)$ l'est -- \deleted{par descente}, on peut pour cela supposer $S=\operatorname{Spec}(k)$, et opérer par descente \ldots Ici il faut $S$ \deleted{\unsure{noe.}} \inserted{loc.} noeth.}

\margin{\uline{Question} Si $T/S$ est fidèl'mt plat sans plus, et $\deleted{\unsure{X}}Y/S$ est tel que $Y\times_S T/T$ soit propre, est-il vrai que $Y/S$ l'est ? -- Oui dans le cas $S$ loc. noeth, car on utilise encore le lemme de Chow.}

[N.B. cette réserve doit pouvoir être levée par la technique de descente]. De $=$, \inserted{\unsure{dans le $=$ voisin. et récipro ? réserve|dans le $=$ énoncé et récipr. ? réserve}.} $f_*$ transforme \deleted{\unsure{applicat.|affines}} morphismes affines en morphismes affines [en effet, soit $\varphi:X\unsure{-|\to}Y$ affine, \deleted{alors} on est ramené \inserted{par excision} au cas où $Y$ est affine et $X\subset Y\deleted{\unsure{\times|\otimes}}\times_T Z$ \deleted{\unsure{$Z/T$, $Z$|$Z$, $\subset$, $Z$}} $X$ ss schéma fermé, \inserted{$Z/T$ affine,} et on se ramène au cas où $X=Y\mathbin{\otimes}_T Z$ puis on traite des imm. fermées] il transforme morphismes projectifs en morphismes projectifs si $T/S$ est séparable, et en tous cas il transforme morphismes quasi-projectifs en morphismes quasi-projectifs.

[$\overline{\mathrm M}$ méthode, du moins si \inserted{si on $X_{/T}\to Y_{/T}$} $Y$ est \deleted{\unsure{noeth.}}. \deleted{\unsure{donc}} qu. proj. de type fini sur $T$ donc un faisceau coh. sur $Y$ est quotient d'un provenant de $T$ ; Il vaut mieux, et il est plus précis, de prouver encore que : un

% archive sheet number: 43
